% SOURCE_STATEMENT: "Can three-dimensional Wess–Zumino–Witten models be generalized to produce trigonometric and elliptic R-matrices for broader choices of the defining differential form?"
% Preparation and final self-audit: CONTRIBUTING.md prompts C and D; source/audit artifacts retained in the local physics_statements_2026_10_11 batch.
% !TeX program = lualatex
\documentclass[11pt]{article}
\usepackage[margin=25mm]{geometry}
\usepackage{fontspec}
\setmainfont{Latin Modern Roman}
\usepackage{amsmath,amssymb,mathtools}
\usepackage{unicode-math}
\setmathfont{Latin Modern Math}
\usepackage{xurl,hyperref}
\hypersetup{colorlinks=true,urlcolor=blue}
\setcounter{secnumdepth}{0}
\setlength{\parindent}{0pt}
\setlength{\parskip}{5pt}
\setlength{\emergencystretch}{3em}
\newcommand{\sref}[2]{\hyperlink{source:#1:#2}{[S#2]}}
\newcommand{\cataloguescope}[1]{\paragraph{Recorded target.}#1}
\newcommand{\R}{\mathbb R}
\newcommand{\C}{\mathbb C}
\newcommand{\N}{\mathbb N}
\newcommand{\Z}{\mathbb Z}
\newcommand{\Q}{\mathbb Q}
\newcommand{\F}{\mathbb F}
\newcommand{\eps}{\varepsilon}
\begin{document}
\section{Trigonometric and Elliptic R-Matrices from Three-Dimensional WZW Models}\label{statement}

\subsection{Definitions and mathematical statement}

\paragraph{Shared notation.}$\mathbb N=\{1,2,\ldots\}$; the standard number fields have their usual meanings. All additional conventions are stated locally.


An $R$-matrix acts on a tensor product and obeys a Yang--Baxter compatibility equation. Rational, trigonometric and elliptic refer to its dependence on a spectral parameter. In the report's three-dimensional Wess--Zumino--Witten construction, the chosen differential form is part of the action and controls the resulting integrable structure.

\paragraph{Catalogue wording retained for traceability.}
Can three-dimensional Wess–Zumino–Witten models be generalized to produce trigonometric and elliptic R-matrices for broader choices of the defining differential form?

\paragraph{Scope and source qualification.}The source asks for an extension of its particular construction. It does not assert that every arbitrary differential form produces an integrable theory.
\subsection{Short English statement}
Extend the three-dimensional WZW construction to broader differential-form choices producing trigonometric or elliptic R-matrices.
\subsection{Sources}
\begin{itemize}
\item[\hypertarget{source:statement:1}{S1}] ULAM.AI and the UnsolvedMath Contributors, UnsolvedMath, catalogue entry. Catalogue curation under CC BY 4.0; definitions and English exposition added locally.
\url{https://huggingface.co/datasets/ulamai/UnsolvedMath}.
\item[\hypertarget{source:statement:2}{S2}] Mini-Workshop: Three Facets of R-Matrices (2021), Oberwolfach Reports, DOI 10.4171/owr/2021/51, pp. 2791--2825. Original report and mathematical context.
\url{https://doi.org/10.4171/owr/2021/51}.
\end{itemize}
\label{end:statement}
\end{document}
